\chapter{Research Significance}
The polar regions are losing ice, and their oceans are changing rapidly~\cite{Meredith_2019}. The consequences of this extend to the whole planet and it is crucial for us to understand them to be able to evaluate the timing and magnitude of potential mitigation strategies. 
The accelerating loss of continental ice sheets in both Greenland and Antarctica are major contributors to global sea level rise~\cite{Meredith_2019}. Impacts extend beyond direct ice loss: as fresh water from melting ice sheets is added into the ocean, it increases ocean stratification disrupting global thermohaline circulation~\cite{Jacobs_2004}. In addition, cold freshwater can dissolve larger amounts of $\mathrm{CO_2}$ than regular ocean water creating corrosive conditions for marine life~\cite{Meredith_2019}.
While there is high confidence in current ice loss and retreat observations in many areas, there is more uncertainty about the mechanisms driving these changes and their future progression~\cite{Fox-Kemper_2021}. Uncertainty increases in regions with variable bed conditions, where characteristics like bed slipperiness and roughness have been difficult to verify via direct observations. Other problematic areas involve the ice sheet grounding line (GL): The zone that delineates where the ice sheet is grounded on bedrock. The GL retreat rate depends crucially on topographical features like pinning points~\cite{Fox-Kemper_2021}, which lead to increased retardation of upstream flow. Although this mechanism is established, major knowledge gaps persist in mapping bed topography across Antarctic ice sheet margins, with over half of all margin areas having insufficient data within $5\!\,$km of the grounding zone~\cite{RINGS_2022}. Addressing this data gap through systematic mapping enhanced by auxiliary data streams with more complete coverage, would significantly improve both our understanding of current ice dynamics and the accuracy of ice-sheet models projecting future changes.

 